\subsection{逆tangent非退化分枝の全値coverを閉じる}
\begin{theorem}\label{thm:h3-order2-22-reverse-plus-empty}
逆tangent・非退化正topの正Pfaffian全領域にH3対象は存在しない。
従って負tangent型に残るH3候補は、逆割当ての平方正top
\[
 [x,y]+uv+vu=0,\qquad x^2+u^2=0
\]
というrank2 cross境界だけである。
\end{theorem}
\begin{proof}
$z_5\ne0$ の保存torusで $z_5=1$ とし、全二パラメータmainを
\[
 (z_0,z_1,z_2,z_3,z_4,z_5)
 =(1+\beta t+t^2/4,\beta,-t/2,t/2,t,1)
\]
と書く。元の全三次・四次関係に対する商汎関数を計算する。
その分母は、次の因子以外に零点を持たない：
\[
 \begin{split}
 D={}&t(\beta-1)(\beta+1)(2\beta+t)\\
 &\cdot(2\beta t+t^2+1)(4\beta t+t^2+3)(4\beta t+t^2+4).
 \end{split}
\]
正常行列の指定 $10\times10$ minorは厳密に
\[
 -\frac{(t-1)^2(t+1)^2}
 {2(2\beta t+t^2+1)^2(4\beta t+t^2+3)(4\beta t+t^2+4)}.
\]
従って $D(t-1)(t+1)\ne0$ で正常三次核は零。
$t=1$ は前節で全正常pencilの因数分解により排除済みである。

残る全曲線を、別の一パラメータ $s$ による次の八行で覆う。
各行で元の関係から改めて商汎関数を作り、二つの正常10minorを計算する。
表の最終列 $E(s)$ は、全消去分母と二minorのgcdの零点を合わせ、
既知の $t=1$ とparameterの定義域外を除いた平方因子のない式である。
\[
\begin{array}{c|c|l}
 \beta&t&E(s)\\\hline
 s&0&s\\
 s&-1&(s-1)(s+1)(2s-1)(4s-5)\\
 1&s&s(s+1)(s+2)(s+4)\\
 -1&s&s(s-2)\\
 -s/2&s&s(s^2-3)(s^2-4)\\
 -(s^2+1)/(2s)&s&(s+1)(s^2-2)\\
 -(s^2+3)/(4s)&s&(s+1)(s^2-3)(s^2-9)\\
 -(s^2+4)/(4s)&s&(s+1)(s^2-2)(s^2-4)
\end{array}
\]
最後の三行では $s\ne0$。
各 $E\ne0$ で正常核は零であり、全特殊点は次の有限集合に入る：
\[
\begin{gathered}
 (0,0),(1,-1),(-1,-1),(1/2,-1),(5/4,-1),\quad
 (1,0),(1,-2),(1,-4),\\
 (-1,0),(-1,2),(-1,3),(1,-3),\quad
 (-3\sqrt2/4,\sqrt2),\quad(-\sqrt3/2,\sqrt3)
 \quad\text{及び後二点のGalois共役}.
\end{gathered}
\]
表の各既約例外因子がこの集合へ写ることも、剰余式の恒等零によって検証する。
有理点は $\mathbb Q$ 上、後二組はそれぞれ二次体上の厳密消去を行う。
$(\beta,t)=(1,-1)$ 以外は正常核が零である。
唯一の例外では正常核が一次元で、その全非零元は
\[
 (x-2y)\bigl((x-2y)^2+4v^2\bigr)
\]
に比例するので商は非domainとなる。
ここで商空間の特殊化を一般点のrankだけから推論せず、各有限点を元の関係で再消去している。

$z_5=0,z_0\ne0$ は $x,y$ の交換と $v\mapsto-v$ によりmainへ移る。
この変換は正Pfaffian atlasで $z_0,z_5$ を交換し、$z_1,z_4$ を保つ。
残る $z_0=z_5=0$ では、必ず $t=z_4\ne0$、$z_1=-t^{-1}-t/4$。
正側二軸は負側を
\[
 f=\begin{pmatrix}t/2&t^{-1}-t/4\\-t&t/2\end{pmatrix},\qquad g=f^{-1}
\]
でnormalizeする。偶二関係から負二次中心 $Z_2$ が得られ、補助次数の議論も
補題\ref{lem:h3-order2-22-odd-determinant-centralizer}の証明と同じである。
ただし今回は $\det f=1$ なので、奇次数を分けて確認する。
補助次数 $\pm1$ では $\operatorname{Sym}^2 f=I$ が必要だが、
$t\ne0$ により $f\ne\pm I$ なので不可能。
補助次数 $\pm3$ は $x^3,y^3$ のスカラー倍だけであり、負側との反交換には $f^3=-I$ が必要。
この行列式の成分と $t$ の可逆性を消去すると、正確に $t=1$ が要求される。
それは前節の三次norm恒等式により全て一次因子を持つ。
従ってこの最後の境界にもH3はない。

検証器
\[
\texttt{verify\_cubic\_order2\_22\_reverse\_plus\_cover.py}
\]
は二パラメータminorを分母消去後の多項式行列式として検証し、
八行の元関係・二minor・例外多項式、全有限点と二次体、唯一の因子、
最後の $f^3=-I$ の局所化idealを再生する。
全AS4存在は仮定せず、各H3候補でのHilbert値と包絡domain性に基づく必要排除である。

負tangent第一割当ては既に全Jordan型で空、逆非退化の負Pfaffianも空である。
非退化正topにはJordan部分がないため、以上を合わせて定理の最後の主張が従う。
逆平方型のrank2 crossには、この可逆graph計算を適用していない。
\end{proof}
